\section{Conclusion and Limitation}
\label{sec:conclusion}
We identify \textit{functional generalization} as an important yet underexplored problem, where an agent must accomplish the same function with unseen tools.
We introduce \algoname{}, a two-stage framework that decouples action-free functional reasoning from grounded execution. Across hitting, sweeping, and hooking, simulation and real-world experiments show improved generalization to unseen tools over existing baselines. Further analyses show that \algoname{} benefits from heterogeneous action-free data, including human videos, and can integrate VLM-based functional point selection for more autonomous tool use.

% \textbf{Limitations and Future Work.}
% While \algoname{} focuses on adapting functional motions across unseen tools, general tool use also requires dexterous interaction, including grasp selection, pose adjustment, and in-hand manipulation. Future work could combine keypoint-based functional reasoning with dexterous end-effectors to jointly reason about grasping and tool use. Our current representation is also 2D; extending it to 3D or multi-view keypoint trajectories could improve spatial reasoning and robustness to occlusion and viewpoint changes.
%\textbf{Limitations and Future Work.}
% Our current multi-source action-free pretraining only incorporates human videos in addition to simulation data. Future work could leverage larger-scale egocentric or web video data, covering more diverse tools and interaction behaviors, to further improve functional generalization. In addition, \algoname{} focuses on functional motion reasoning and does not explicitly model dexterous interaction, such as grasp adaptation or in-hand manipulation, which is also essential for general tool use and could be incorporated in future work.
\textbf{Limitations and Future Work.}
\algoname{} remains limited in generalizing to unseen tools that require substantially different functional motions. Future work could leverage larger-scale egocentric or web videos with more diverse tool-use behaviors to further improve functional generalization. In addition, our method focuses on functional motion reasoning and does not explicitly model dexterous interaction, such as grasp adaptation or in-hand manipulation, which is also essential for general tool use and could be incorporated in future work.